\documentclass[acmsmall,screen]{acmart}

%% ---- FSE 2026 / PACMSE (Proceedings of the ACM on Software Engineering) ----
%% Named (non-anonymous) version of the focused paper (source: fse_focused_v2.md).
%% v3: adds Section 9, a separately and later pre-registered live study that tests
%% the detection-cost coupling on a second model family (GPT-5.5). No change to the
%% original V1->V2 narrative or its frozen numbers; see the Integrity log (D-V3-1).
\settopmatter{printacmref=false, printccs=false}
\setcopyright{none}
\acmJournal{PACMSE}
\acmVolume{1}
\acmNumber{FSE}
\acmArticle{1}
\acmYear{2026}
\acmMonth{7}

\usepackage{booktabs}
\usepackage{tabularx}
\usepackage{enumitem}
\usepackage{listings}

\lstset{basicstyle=\ttfamily\footnotesize, breaklines=true,
  breakatwhitespace=false, columns=fullflexible, keepspaces=true,
  xleftmargin=4pt, aboveskip=6pt, belowskip=6pt, frame=none,
  extendedchars=true, inputencoding=utf8, literate={§}{{\S}}1}

\begin{document}

\title{When the Monitor Blinds Itself: Failure Detection in Multi-Agent LLM
Systems Is Limited by Observation---and Coupled to Its Cost}

\author{Anshuman Mullick}
\affiliation{%
  \institution{Bilkent University}
  \city{Ankara}
  \country{Turkey}
}
\email{anshumanmullick7@gmail.com}

\author{Eray T\"uz\"un}
\affiliation{%
  \institution{Bilkent University}
  \city{Ankara}
  \country{Turkey}
}
\email{eraytuzun@cs.bilkent.edu.tr}

\begin{abstract}
A multi-agent LLM system runs many workers in parallel against a single plan.
When that plan quietly breaks mid-run---a renamed dependency, a removed
endpoint, a drifted schema---the workers keep spending tokens on a dead plan
until the very end. We built a monitor to catch this early. We wrote down, in
advance, the bar it had to clear, and we ran it to a verdict we were not free
to renegotiate. It failed, and the failure is the finding.

The finding has two halves. First, catching a broken plan is limited by what
the monitor gets to \emph{see}, not by how clever it is. A correct check
catches nothing on a part of the world that no one looks at again after it
breaks. And a monitor that interrupts too eagerly spends the run's budget on
false alarms---the same budget a second look would have needed---so its
accuracy and its restraint draw on one finite pot.

Second, the only mechanism that restores that second look---re-reading the plan
and writing fresh checks for it---is also the system's biggest cost. So cutting
the cost risks the very detection it pays for: \textbf{detection and cost are
coupled.}

We establish both halves on our own system. Our first version failed badly:
35\% detection, a ``judge'' stage that approved all 18 of its own false alarms,
and a cost worse than the plain baseline it was meant to save. A \emph{naive}
baseline beat it at zero false alarms. A trace-level diagnosis, checked by two
outside red teams, found the cause: the monitor starved itself of
observation---on identical seeds, ordinary baselines re-looked at every starved
surface (12 of 12) while ours did not. The redesign restores observation
directly. It catches 24 of 31 injected faults at zero false alarms on clean
runs, and it even transfers to a fault type it was never shown (3 of 5, where
every passive baseline scores 0 of 5). But it still misses its 12\% cost bar at
55.5\%: an \textbf{overall fail} on cost. A cost autopsy traces nearly all of
that overhead to one step---re-reading the plan to write the checks---which is
exactly the step that produces the detection. The obvious escape, spreading
that one-time cost across more workers, does not open: in a scaling pilot the
cost gap \emph{widens} as we add workers. Restoring observation fixes detection
and ends the false-alarm self-harm. It does not yet pay for itself.

Finally, we asked whether the coupling is one model family's quirk. In a
\emph{separately and later} pre-registered live study, we swapped a different
vendor's model (GPT-5.5) into the check-writing step through a frozen seam. It is
a cheaper, higher-quality writer---yet the frozen bar failed a third time.
Offline the escape looked open (detection is flat as we cut the model's
reasoning); live it closed, because the cost lever that could clear the bar is
\emph{checklist coverage}, and coverage is what does the catching. Shortening the
checks to save money drops exactly the faults that need a specific probe---live
detection falls 5/5 to 4/5 to 3/5---and overhead never reaches 12\%. The coupling
holds on a second model family, through a new mechanism.
\end{abstract}

\keywords{multi-agent LLM systems, runtime verification, self-adaptive
systems, failure detection, interrupt policy, pre-registration, agent
orchestration}

\maketitle

\section{Introduction}
Twelve seconds before the fault it existed to catch, our monitoring system
destroyed its own coverage.

Here is what happened in one run. A check fired on a perfectly healthy file
listing. A ``judge'' model approved the alarm as real, at 0.98 confidence, by
hallucinating a misreading of the listing. The orchestrator paused its workers
and replanned. The new plan got a fresh set of checks---and those checks
covered none of the surfaces the old ones had watched. Twelve seconds later the
real fault fired into a world no one was watching. The run finished, confident
and wrong.

Every stage did its job. The system as a whole blinded itself.

This paper is about that mechanism. In a multi-agent system on a fixed budget,
catching a broken plan is bounded by \emph{observation}, not by the monitor's
\emph{intelligence}---and a monitor that is too eager starves its own
observation. We establish this through three acts: a pre-registered experiment
that killed most of our design, a careful diagnosis of why, and a redesign that
puts observation back at the center.

\textbf{The stakes.} Multi-agent LLM systems are expensive. They use on the
order of 15$\times$ the tokens of single-agent chat \cite{ref27}. When a plan
breaks mid-run, plain batch orchestration only finds out at the end---after
every parallel worker has spent that multiplier on a dead plan. These failures
are real and well-catalogued \cite{ref14, ref16}. And the \emph{mechanism} to
interrupt a run already exists in production frameworks \cite{ref22}. What is
missing is the \emph{policy}: something that decides what to watch for, whether
an anomaly actually breaks the plan, and when interrupting is worth it.

\textbf{What we built.} We proposed a layer called Sentinel Protocol. The idea
is simple. Before the workers run, an LLM reads the plan once and writes out a
set of small, concrete checks---we call them \emph{tripwires}---that say what
each step is trusting about the world. The workers then watch for those checks
cheaply, by pattern matching, and a confirmed break interrupts the orchestrator
so it can replan. We pay for expensive reasoning once, up front, and keep the
running system cheap. (For software-engineering readers: it is a classic
self-adaptive control loop \cite{ref23}, but its monitor is \emph{generated
from the plan at runtime} rather than written by hand.)

\textbf{What we did to it.} We then did something the field rarely does to its
own proposals. We wrote down pass/fail lines \emph{in advance}---with the
consequence of each failure decided ahead of time---froze them, and ran the
experiment to a verdict we could not move after seeing the data.

The verdict killed most of the design. Detection was 35\%, against a 60\% bar
and a 40\% floor below which the claim died outright. The ``judge'' stage
approved all 34 alarms it ever saw, including all 18 false ones---a rubber
stamp, not a filter. The cost was worse than the plain baseline it was built to
save. One claim survived: where a fault keeps showing up on a surface someone
watches, our checks catch it about 3$\times$ faster than a cost-matched
periodic re-check (median 3 vs.\ 9 tool calls).

\textbf{Why it failed.} A failed verdict is only useful if you find out why. We
replayed the fully recorded runs, sent the diagnosis to two outside red teams
for blind review, and settled every testable claim with a deterministic replay
that cost no LLM calls. The diagnosis: detection here is observation-bounded,
and in our case the bound was self-inflicted. On identical seeds, ordinary
baseline systems re-looked at every starved surface (12 of 12)---but ours did
not. Its own false alarms had eaten the part of the run in which a second look
would have happened. Under a fixed budget, accuracy and restraint are not
separate dials.

\textbf{The redesign.} The fix follows from the diagnosis: stop relying on
workers passively noticing things, and instead have the monitor \emph{re-read
the world itself} with cheap, deterministic checks. We validated this mechanism
in offline replay before building it. Then we re-measured the new version
against a fresh set of pre-registered bars, frozen before any new data existed,
with two fault types whose exact settings were sealed away with someone outside
the project---so the design could not have seen them.

\textbf{The new verdict.} The redesign passed detection and failed cost. It
caught 24 of 31 injected faults, at \emph{zero} false alarms on clean runs,
recovering exactly the quiet, hard-to-see fault types the first version was
blind to. It even transferred partway to the sealed fault types it had never
seen: it caught the held-out RESOURCE\_BUDGET type 3 of 5 times, where every
passive baseline caught it 0 of 5. (The other held-out type,
DEPENDENCY\_VERSION, was missed by \emph{every} system, ours included---a limit
of the benchmark, not of our design, as \S7 explains.) But its extra cost on
clean runs was 55.5\%, against a 12\% cap---an \textbf{overall fail}, driven
entirely by cost.

So the thesis sharpens. Restoring observation fixes detection and ends the
self-harm. But the restored observation also has to be \emph{cheap}, and ours
is not.

\textbf{This is not a system that failed at its only job.} Where signals
recur, it detects about 3$\times$ faster than the obvious alternative. It
injures \emph{zero} clean runs, where both its own predecessor and the naive
baseline raise false alarms. And it carries part of its skill to fault types it
was never shown. The cost failure is not the absence of value. It is that the
value does not yet clear the cost of looking. Section 8 shows that this cost
failure and the detection success are two views of the same fact. Section 9 then
shows the same fact on a different vendor's model, under a separate frozen bar
committed later: a cheaper, better writer, and the coupling holds anyway.

\textbf{Contributions.}
\begin{itemize}
  \item \textbf{The finding} (\S6, \S8, \S9): in budget-bounded multi-agent
  execution, detection is limited by observation, that limit can be
  self-inflicted, and---because the one thing that restores observation is also
  the main cost---detection and cost are \emph{coupled}. The negatives are not
  a string of missed targets; they converge on one fact. That fact reproduces on
  a second model family (\S9): the coupling is not one vendor's artifact.
  \item \textbf{A pre-registration method for agent-systems research} (\S5,
  \S7, \S9) that actually bound us: frozen bars returned \emph{fail} \emph{three}
  times---twice on both versions of our own system, and a third time in a
  separately and later pre-registered live study of a different vendor's model
  (\S9)---and the pre-decided consequences ran as written each time. We add
  binding outside red-team review, a data embargo, a log of every run, sealed
  held-out fault types, and a check-writer that is never told the fault list.
  \item \textbf{A diagnosis method} (\S6): deterministic replay over fully
  recorded runs, used to settle outside criticism with evidence rather than
  argument, and to test the redesign before building it.
  \item \textbf{TripwireBench} (\S4): a benchmark of injected faults that are
  \emph{verified to actually break a plain baseline}, with clear recovery
  labels, sealed held-out types, and cost/speed/noise metrics.
  \item \textbf{The surviving result and the naive baseline} (\S5, \S7): a
  3$\times$ speed edge where signals recur, and a no-frills baseline that
  caught 56\% at zero false alarms and beat our full first version in every
  category---the honest floor any successor has to clear.
\end{itemize}

\section{Background and Related Work}
\textbf{Self-adaptive systems.} Software-engineering readers will recognize the
shape: a monitor--analyze--plan--execute control loop, a staple of
self-adaptive systems for two decades \cite{ref23, ref24}, including recent work
that puts LLMs inside the loop \cite{ref25, ref26}. We do not claim the loop.
What is new is where its parts come from: the monitoring logic is \emph{written
from the plan, at runtime, by an LLM} rather than designed by hand. Our
experiment then shows that the analysis step is the hard part---an uncalibrated
filter does not just fail to help; it does damage (\S5).

\textbf{The nearest prior work, confronted.} The closest ancestor is RVPLAN
\cite{ref6, ref34}, which turns a planner's stated preconditions into runtime
monitors and replans when one breaks. ``Watch the plan's assumptions, interrupt
to replan'' is therefore established. But RVPLAN lives in a setting where three
things hold that do not hold for us:

\begin{itemize}
  \item \emph{The specification already exists.} RVPLAN translates
  preconditions that sit, machine-readable, in a planning file. Our checks have
  to be \emph{invented} from a plan written in natural language. Coverage is
  therefore bounded by what the writer can infer (\S6).
  \item \emph{The world narrates itself.} RVPLAN reads a clean stream of
  labeled facts in its own vocabulary. We read raw, unlabeled tool output,
  often only once---and the gap between ``watching'' and ``actually seeing'' is
  our central result.
  \item \emph{The stream is clean.} In a labeled stream a precondition simply
  holds or it does not; nothing can false-fire, so nothing needs filtering. In
  our setting the filtering layer is the whole game, and it is where our first
  version broke.
\end{itemize}

An established paradigm meets the LLM-agent setting and breaks in three
measurable ways. This paper is about what it takes to survive there.

\textbf{Other lines.} Runtime verification and monitor synthesis \cite{ref4,
ref5} give us the spec-to-monitor tradition; we trade full temporal logic for a
small predicate language over tool calls, and our complaint is about
\emph{passivity}, not the trade (\S3). LLM-agent guardrails \cite{ref9, ref10,
ref11, ref12, ref13} enforce externally written \emph{safety} rules; our checks
come from the plan, target plan validity and cost, and feed replanning. Failure
taxonomies \cite{ref14, ref20} catalog things agents do wrong; we cover the
complement---things the \emph{world} does that the plan didn't expect---and we
use them to write checks before the run, not to label failures after.
Observability tools \cite{ref15, ref16, ref17, ref18} explain failures after
the fact, for humans; we act on the same signals in flight. LangGraph's
interrupt \cite{ref22} is mechanism without policy; schedulers \cite{ref19}
consume interrupts once someone decides to raise them. We sit in the gap:
deciding \emph{what to watch} and \emph{whether to pull the trigger}.

\section{The System Under Test}
This section describes the system, but the system is our \emph{instrument}, not
our contribution. It is the thing we pre-registered, killed, diagnosed, and
rebuilt in order to reach the finding. We describe the current (v2) version and
note where the first version (v1) differed. Every change from v1 to v2 cites
the trace evidence that forced it (\S6), and v1's verdict is reported in full
no matter how v2 turns out (\S5, \S7).

\textbf{Three roles.} The \emph{orchestrator} owns the plan and is the only
part allowed to replan; it acts only on confirmed breaks. The \emph{sentinel}
runs once, up front: it reads the plan and writes the checks. The \emph{workers}
carry those checks and watch for them cheaply while they work, raising a
structured alarm on a match. The cycle is: plan, write checks, run, raise an
alarm, confirm-and-interrupt, replan. Pay for the thinking once; keep the
running system cheap and deterministic. The pilot added one rule the new
version follows: the running system also has to \emph{look}---not just wait to
be told.

\textbf{Checks and probes.} A check is concrete. It names a specific thing to
watch (a status code, a field, a structure, a hash), a specific value, and what
to do if it trips. Vague checks (``if the API changed'') are not allowed. One
rule we learned the hard way and now enforce: every constraint the check-writer
must follow has to live in a place the model actually reads---a field
description in the schema---because rules written in comments or prose simply do
not reach it.

The pilot showed the checks were fine but \emph{passive}---they only noticed
what a worker happened to read. So the new version's key addition is the
\textbf{active probe}. Each check can carry a small, deterministic re-read of
the world (a HEAD request, a schema fingerprint, a status re-check) that runs on
its own, with no LLM involved. Three things, each forced by what we measured
(\S6), shape how probes run:

\begin{itemize}
  \item They run on a \emph{separate channel} that does not disturb the world
  they measure---replay found three concrete ways a naive probe perturbs its
  own measurement.
  \item They fire \emph{on events and risk}, not on a fixed clock---a naive
  fixed schedule once cost 2,322 probe calls on a 107-call run.
  \item A non-obvious alarm interrupts \emph{only with a confirming probe}. The
  first version's ``wait for a second signal'' rule failed because correlated
  noise confirms itself---6 of 18 false alarms passed it. (A clear status-code
  error still interrupts on its own; no false alarm ever carried a status
  $\ge$ 400.)
\end{itemize}

Probes that compare a value need something clean to compare against. So at the
very start of a run, before any worker acts, the system takes one clean reading
of every surface it will watch, and later compares against that (deviation D30).

\textbf{We removed the judge.} The first version put an LLM ``judge'' between an
alarm and an interrupt. Its record: it approved all 34 alarms it ever saw,
ruled the only two genuine signals it down-ruled as noise, and once approved a
hallucinated misreading at 0.98 confidence. We committed---in writing,
\emph{before} the diagnosis results were readable---to making the no-judge
version the primary design (deviation D33 keeps a rebuilt judge only as an
exploratory side arm). There is no after-the-fact selection here.

\textbf{The check-writer is never told the fault list.} The first version
walked a list of eight named fault categories. The new version's check-writer
gets \emph{no list}. It reasons from two general ideas instead: what does this
step trust about the world, and which of six general \emph{shapes of change}
might break it (a field vanished, a status moved, a structure changed, a value
moved, an order scrambled, a relationship broke). Its only worked examples come
from the five fault types we had already seen, chosen by a rule fixed before any
prompt tuning (deviation D27). We attach category labels \emph{after} detection.
This is what makes the held-out test honest: the writer has to \emph{transfer}
general reasoning to a sealed fault type, not look it up (\S7). The pilot's
sharpest lesson is why this matters. Fault types differ enormously in how
visible they are. The one type the first version caught well was the loud,
repeating, status-coded one; the quiet, seen-once types were nearly invisible.
The probe exists to make every type as visible as the loud one.

\section{TripwireBench}
To test a monitor, you need a world that breaks on cue.

TripwireBench is that world: a fully mocked, deterministic set of services
(auth, inventory, pricing, shipping, a document store, a code repo), with
seed-derived data and no clock values anywhere. Faults are exact state changes
that fire at exact tool-call counts, so any run can be replayed byte-for-byte
from its setup. Because a fault fires at the same \emph{point in the work}
across systems and seeds, batch and treatment runs are directly comparable.
Four task types (API aggregation, repo migration, research synthesis, document
pipeline) run on a real agent CLI; one piece of middleware counts calls, fires
faults, checks tripwires, and writes a full trace.

\textbf{Every fault has to actually hurt.} No fault enters the experiment until
it provably breaks the plain batch baseline: injected runs must fail while clean
runs pass, under a rule fixed in advance. The pilot qualified all 9 fault--task
pairs. Two more fault types were authored \emph{after} the pilot and held out: a
metered quota that only complains after the budget is already spent
(RESOURCE\_BUDGET), and a silent version bump that quietly truncates a default
page of results (DEPENDENCY\_VERSION). Their exact parameters were sealed with
someone outside the project, so part of the scoring could not be over-fit. Both
cleared the same bar---the quota fault on the first try, the version fault only
on its third revision, with the two failed designs reported as findings rather
than hidden. The qualification record is itself instructive: a competent batch
planner kept \emph{defeating} faults---by paginating defensively before it saw
anything, or by reading collections once at the start so a later change landed
after the only look---and \emph{strengthening} the baseline un-qualified two
faults that had qualified earlier, so wounds must be re-checked whenever the
baseline improves (deviation D23).

\textbf{The systems we compare.} S1, plain batch (dispatch, wait, aggregate,
one redo). S2, naive interrupt (workers raise any anomaly straight to the
orchestrator---no check-writer, no judge). S3, a cost-matched periodic
re-check, tuned so its overhead equals our monitor's---the strongest fair
version of the obvious alternative. S4, our monitor without the judge. S5, our
full first version. The confirmatory study runs four arms---S1, S2, S3, and the
redesigned monitor---setting the unbuilt rebuilt-judge arm aside (deviation
D33).

\textbf{What we measure.} Wasted work (tokens and dollars spent after a fault
that don't reach the final answer); detection speed (calls from fault to first
pause); detection rate; false-alarm rate (with no fault and no alarm scored as
zero); extra cost on clean runs; and plain pass/fail task success---no LLM
grading anywhere. Pause attribution follows a written rule, with a hand-audited
sample (13 of 13 agreed).

\textbf{Integrity machinery.} The benchmark's credibility tooling is part of
the artifact: a frozen pre-registration with cryptographic fingerprints, a
numbered log of every departure, fault triggers that only the test harness can
reach, a log of every run touching any treatment component, a data embargo
before the bars were frozen, and sealed held-out parameters verifiable against a
public hash.

\section{The Pre-Registered Pilot and Its Verdict}
\textbf{Setup.} 13 task variants (9 qualified injected pairs + 4 clean) $\times$
5 systems $\times$ 3 seeds = 195 runs, all banked, with no malformed traces.
Tooling was pinned, model assignments fixed per role, sessions isolated. Faults
fire at the halfway point of each task's typical length, identical across
systems and seeds. Before the main experiment, two independent outside models
scored the written checks against a four-part rubric (95.1\% agreement) and they
passed at 89\% coverage against an 80\% bar. The four pass/fail bars were frozen
before the first run, each with its consequence decided in advance. The decision
window after the verdict was itself pre-set: 48 hours, ``do not negotiate with
the data''.

\textbf{The verdict (computed once, shown verbatim):}

\begin{lstlisting}
KG1 recall: 35% (>=60%) | categories >=50%: 1/5 (>=4)            -> FAIL
KG2 FIR: S5=1.0 S2=0.0 S4=1.0 (S5<=0.5*S2 and S5<=0.7*S4)        -> FAIL
KG3 cost: S5 med $1.178952 vs S1 med $0.340831; success S5=4%
     S1=22%; clean overhead OVER (<=12%)                          -> FAIL
KG4 vs heartbeat: wasted med S5=15014.0 S3=7950 (>=20% better)
     | TTD med S5=3 S3=9 (>=2x)                                   -> PASS (TTD arm)
\end{lstlisting}

Plainly: the bars we froze returned \emph{fail} on our own flagship, and the
pre-decided consequences ran as written. Detection at 35\% fell in the kill
band (below 40\%), so that claim died. The false-alarm gate failed, so we
dropped the judge. The cost gate failed, which triggered a break-even fit that
later killed the cost claim outright (\S8). Only one thing survived: speed. Our
checks caught the fault in 3 calls to the periodic re-check's 9. But the
wasted-work side of that same gate failed (15,014 vs.\ 7,950 tokens):
\textbf{speed did not turn into savings under the first version's machinery}.

\textbf{Reading it carefully.} ``Detection'' here means a confirmed,
trace-attributable pause, defined before the bars computed. A more generous
reading (counting alarms that didn't pause) recovers only two runs---so the low
number is real silence, not a strict definition. By category: the one passing
type (PERMISSION\_AUTH, 5/6) is exactly the loud, repeating, status-coded one;
the quiet types scored 2/6, 1/6, 1/6, and 0/3. The predicted near-zeros showed
up.

\textbf{The naive baseline, stated plainly and prominently.} S2---workers raise
any anomaly straight to the orchestrator, no check-writer, no judge---caught
\textbf{15 of 27 (56\%) at a false-alarm rate of zero}, and beat or tied our
full first version in \emph{every} category. Workers given no filter escalate
rarely, and mostly correctly, at this scale. The folk worry that naive
escalation means alarm spam---which our false-alarm gate had baked in---was
simply false here. Sophistication actively hurt. S2 is now a required
head-to-head arm, under a rule we committed to in advance: if the simple thing
wins, we say so in the results, not a footnote.

\textbf{It hurt healthy runs, too.} On clean runs with nothing wrong, the full
first version succeeded 3 of 12 times against batch's 9 of 12. Seven of the nine
clean failures were self-inflicted: five runs ground to death adjudicating
noise, and two replanned so early that there was nothing to carry forward. All
34 interrupts had been judge-approved; 18 matched no fault at all, and 3 fired
on entirely clean runs. An uncalibrated monitor did not just fail to help. It
did harm.

\section{Why It Failed}
\textbf{Where the misses happened.} We classified all 27 injected treatment
runs along the chain \emph{write check $\rightarrow$ fault fires $\rightarrow$
alarm $\rightarrow$ judge $\rightarrow$ pause $\rightarrow$ attribute}. Of 18
misses: 1 had no check written; \textbf{12 had a check that simply never
fired---the dominant case}; 2 fired but never raised an alarm; 2 were ruled
noise; 1 fault never fired. So the problem was not blindness at write
time---coverage was near-complete once a measurement artifact was fixed. It was
\emph{silence at look time}.

\textbf{We did not grade our own homework.} We sent the diagnosis and the
redesign to two outside red teams for blind review, then settled every testable
claim with a deterministic replay over the recorded runs---with our commitments
fixed before any replay result was readable, and at zero LLM cost. All 27
injected worlds replay identically, which is what makes this possible. The
reviewers split roughly evenly on the eighteen claims they raised, which is
exactly why we \emph{adjudicated} their review with evidence instead of
deferring to it. The decisive checks: the story that evidence reached the system
and was lost is refuted (9 of 12 starved runs had \emph{no} post-fault reading
at all---nothing to lose); the story that false alarms replanned the surface
away is partial (that exact signature appears 0 of 12 times, though 7 of 12 died
of noise); and the story that runs simply died before reaching the fault is
refuted (0 of 17).

\textbf{The load-bearing result.} Our first diagnosis blamed the world: some
surfaces are only ever read once, so nobody could re-look. The same-seed
analysis corrected that, and it is the strongest evidence for the ``not about
intelligence'' half of the finding.

On identical seeds, the ordinary baseline systems re-looked at every starved
surface---12 of 12---and our ``smarter'' monitor did not.

The dumb baselines saw what the smart monitor missed. What starved it was its
own noise: every false alarm spends turns and budget that detection itself
needed. One run logged 172 alarms on healthy surfaces before grinding to death.
One mechanism, four symptoms: noise starved detection, the over-trusting judge
made the noise, and the grinding wrecked clean runs and cost. The opening story
(\S1) is this in a single run.

\textbf{Errata, kept visible.} Two early claims were corrected by the replay and
stay in the record with their corrections: the ``read-once surface'' story was
wrong for its own example (the surface was read \emph{after} the fault, so the
original telling was backwards), and ``one matcher defect'' became a \emph{class}
of dead checks---84 checks that could never match any real path under the live
dialect, fixed with a compile-time liveness test (deviation D24).

\textbf{We tested the fix before building it.} We mechanically built probes for
all nine pilot faults and replayed them against the broken worlds. Probe-first
confirmation blocked \textbf{all 18 false alarms}, and projected detection
\emph{rose} to 21 of 27, because probes recover what alarms lose. We are careful
here: \textbf{every number in this paragraph is a ceiling under stated
assumptions}---hand-built probes the real check-writer had not yet produced,
fixed run paths, the same nine known faults. It is a feasibility test of the
mechanism, and we quote it nowhere as a result.

\section{The Redesign, Measured}
The confirmatory study is the rerun path we wrote down \emph{before} the pilot
verdict existed: report the first version as it stands, diagnose from traces,
change the design with cited evidence, and re-measure against fresh bars.

The integrity rules are in ink. Every shared threshold is inherited from the
\emph{pre-verdict} values---60\% detection, the same category and floor rules,
the 2$\times$ speed bar---so no bar was set by anything the diagnosis told us.
New rules add a hard ``probes must be valid'' gate, a recovery-quality clause
(detect-only, detect-and-recover, and detect-and-justified-stop are counted
separately, so stopping cannot be dressed up as success), and \emph{absolute}
noise caps in place of the broken ratio gate. Fault settings were redrawn from
fresh seeds the designers never saw. The two sealed fault types sit inside the
scoring. A data embargo and a full run log stood throughout. And because the
check-writer never sees the fault list, the sealed types are unseen even in the
vocabulary it reasons with.

\textbf{Result.} The experiment ran to completion---\textbf{172 of 172 runs
banked}, across four systems (S1, S2, S3, and the redesign)---and the bars
computed once.

\begin{lstlisting}
1bKG1 detection: recall 10/15 = 66.7% (>=60%, floor clear);
     5/5 seen categories >=50% (API 6/6, SCHEMA_DRIFT 3/3, AUTH 6/6,
     TOOL_CONTRACT 3/3, RETRIEVAL_INTEGRITY 3/3); probe-validity 10/10;
     recovery-quality 17/24 = 70.8%             -> detection sub-terms PASS
1bKG2 noise/self-harm: clean median FIR 0, P95 0, max 0; 0 grinds;
     pre-detection median FIR 0; clean success 8/12 = 66.7%        -> PASS
1bKG3 cost: V2 clean median $0.3642 vs batch $0.2342;
     clean overhead 55.5% (cap <=12%)                             -> FAIL
1bKG4 vs heartbeat: V2 wasted 7008 vs S3 6404 = 1.09x;
     waste-parity not met (S3 detects 0/31)                       -> FAIL
OVERALL                                                           -> FAIL
\end{lstlisting}

The verdict is \textbf{fail, driven by cost}---strong detection that misses the
cost bar.

One reporting note we keep visible. An admissibility check (a byte-identity
replay) came back 123 of 124 identical; the single difference is a harmless
ordering quirk on the non-detecting periodic-re-check arm, characterized and
logged as deviation D34. The gate \emph{code} folds that check into the
detection gate, so the combined gate prints \texttt{FAIL} even though all four
detection sub-terms pass. We report it split out---detection passes; the
admissibility check is fine bar one benign diff---and note the overall fail is
unaffected, because it rests on cost.

\textbf{Detection, in plain numbers.} The redesign caught \textbf{24 of 31}
injected runs. The passive baselines caught \textbf{0 of 31} (both batch and the
periodic re-check), and the naive baseline caught \textbf{12 of 31}. On the five
seen fault types it was perfect, and it finally caught the quiet, seen-once
types the first version could not (\S5). Two numbers describe the same runs from
different angles: 24 of 31 \emph{detected} (77\%), and 10 of 15 (66.7\%) on the
gated subset---the subset that admits a recovery path. They differ by
\emph{definition}, not contradiction. The 24 detections split into 3
detect-and-recover, 14 detect-and-justified-stop, and 7 detect-only, with 17 of
24 (70.8\%) in the two passing buckets. Per-category numbers carry Wilson 95\%
lower bounds; three categories sit at n=3, where one run swings the rate 33
points (reported with the bound).

\textbf{It transferred to a fault type it had never seen.} The redesign caught
the sealed RESOURCE\_BUDGET type \textbf{3 of 5} times, where every baseline
caught it \textbf{0 of 5}---real transfer of general reasoning to a fault type
held out of its training entirely.

The second sealed type, DEPENDENCY\_VERSION (a silent version bump that
truncates a default page), was missed by \emph{every} system, ours included
(\textbf{0 of 5}). Read-only trace checking over all 20 of those runs explains
why. The fault fires on the very first tool call---\emph{before} the redesign
takes its clean opening reading---so even our reference reading is already
truncated, and the later re-reads match it. And no surface anywhere carries a
version field, so the passive and naive baselines have no error to flag either.
Because \emph{every} system missed it---including the anomaly-raising naive
one---this is a limit of the benchmark, not of our design: it is the same
observation bound from \S6, seen on the read side (a fault that fires before the
monitor takes its first reading cannot be caught by comparison).

\textbf{The first version's pathology, reversed.} The first version's
signature---a rubber-stamped flood of false alarms that starved detection and
hurt clean runs---is gone. Clean-run false-alarm rate is 0 (median, 95th
percentile, and max), with zero grind-deaths. And the win holds under real
pressure: the clean set includes the legitimate config-migration write that used
to trip a false alarm, and the redesign ran it at zero interrupts---while the
naive baseline raised false alarms on four clean runs. The redesign bought back
detection \emph{and} dropped noise to the naive baseline's zero floor. Against
the naive baseline's 12 of 31, sophistication helped this time. But the help did
not survive the cost gate---which is the subject of the next section.

\section{The Coupling: One Result, Not Three Failures}
The cost failures of the first and second versions are not two missed targets.
They are one fact, seen from a few sides: \textbf{detection and the cost of
restoring observation are tied together}, because the one thing that restores
observation is also the main cost.

\textbf{The cost is the capability.} We broke the 55.5\% clean overhead down
call by call. Almost all of it is one new component: the once-per-run LLM step
that reads the plan and writes the checks (median \$0.1376 on clean runs,
against a \$0.130 gap over batch), present in all 12 clean runs. The whole
detection machinery---the probes, the re-readings, the confirmation
logic---costs essentially nothing measured. So the overhead lives in the one
step we cannot remove without removing the capability: that step \emph{is} the
reasoning that produces both the detection and the transfer to unseen fault
types (\S7). Making it cheaper---caching it, shrinking it, skipping it---is not a
free optimization. It risks the detection it pays for. That mutual dependence is
the coupling. A no-check-writer variant makes the point from the other side:
drop that step and the system is nearly cost-neutral with batch on clean
runs---so the step really is the cost center, though that variant's detection is
only established on its narrow home case, so we report it as cost isolation, not
a free lunch.

\textbf{The obvious fix doesn't open.} A fixed per-run cost should shrink, per
worker, as you add more workers. So the natural escape is to spread that
one-time check-writing cost across a bigger fan-out. The break-even model even
predicts a crossover for the redesign---but that is a model prediction, not a
measurement (Appendix B).

We built a width-scaling benchmark and ran a sizing pilot to test it. The
pilot's point estimate---kept out of the formal sizing by design---is that the
redesign is \textbf{costlier than batch at every width and fault rate we tried,
and the gap \emph{grows} as we add workers}, because the check-writing scales
with the planning the system does, not with the parallel work it could dilute
against.

We treat this as a \emph{closing negative}, not a confirmation: the escape shows
no sign of opening, and nailing it down at the pre-registered precision would
cost about ten times the study budget ($\approx$ \$4,430 against a \$450 budget)
for a number already pointing the wrong way---so a full fan-out study is not
warranted. The price tag is the side note. The \emph{direction} is the finding.

\textbf{And it fails even when looking is free.} Our benchmark is a
deterministic mock, where re-reading the world costs nothing. That is the
\emph{most generous} setting the system could ask for---it charges nothing for
the very looking the system exists to do. And it \emph{still} misses the budget,
on the one-time check-writing step alone. So the measured 55.5\% is a
\emph{floor}: in a real deployment, probes hit real endpoints at real cost, and
the volume that reads as free here comes back as real overhead. A system that
cannot pay for itself where looking is free will not pay for itself where
looking costs money.

\textbf{Putting it together.} Restoring observation fixes detection (\S7) and
ends the false-alarm self-harm (\S7). But the only place to cut cost is the
capability-critical check-writing step, the obvious amortization does not open,
and the failure holds even with looking priced at zero. So in budget-bounded
multi-agent execution, \textbf{detection and cost are coupled, and this kind of
monitoring does not yet pay for itself.} That is the next problem to solve---and
it is a fresh experiment, not a patch.

\section{The Coupling Holds on a Second Model Family (Live)}
Everything above is one model family: a Claude check-writer, Claude and Haiku
workers. A fair objection is that the coupling is that family's artifact. So we
tested it on a second one.

\textbf{This study was pre-registered separately, and later.} After the original
pre-registration and both its verdicts, on 2026-07-01, we froze a new, standalone
pre-registration for this study alone---its own 12\% bar, its own arms, its own
pre-decided consequences, its own hard spend cap, all committed before any of its
data existed. It is not part of the original plan and we do not present it as one:
it is an additional, later-dated extension that reruns the paper's method on a
different vendor's model. Folding it into the original frozen pre-registration
would break the exact discipline this paper is about, so we keep the two apart
and dated.

We swap GPT-5.5 (pinned \texttt{gpt-5.5-2026-04-23}; \$5.00 / \$0.50 / \$30.00 per
million input / cached-input / output tokens, verified 2026-07-01 \cite{ref38})
into the check-writing step through the same frozen seam the redesign already
exposes. Everything downstream---grounding, probes, confirmation, scoring---is
byte-identical; only the writer changes.

\textbf{Offline, the escape looked open.} GPT-5.5 is a cheaper and cleaner writer
than the incumbent. It writes valid checklists at \$0.115 per call against the
Claude writer's \$0.1376 (\S8), and it emits the value-comparing check that the
cheap alternative never does (median 8 such checks per checklist, against 0 for a
budget model). We first scored its checklists in the zero-LLM deterministic
replay of \S6. There, detection is \emph{flat} as we turn its reasoning effort
down---\texttt{none}, \texttt{low}, and \texttt{medium} all catch the same 16 of
18 seen faults---because the reasoning tokens are not what does the detecting. At
\texttt{effort=none} it costs \$0.064 a call, well under the incumbent, detection
intact. On the reproducible replay, the coupling looked beaten.

\textbf{Live, it was not.} The replay prices looking at zero and holds the world
byte-identical; it cannot see what a shorter checklist costs in a real run. So we
ran the priced matrix live---real workers, real timing---and measured the
overhead the bar is written on. The cost lever that could actually clear 12\% is
not reasoning effort (which is detection-free) but \emph{checklist coverage}: a
shorter checklist is a cheaper writer call. We capped it with a single general
instruction, committed in the pre-registration and blind to the fault list---keep
the $N$ most load-bearing assumptions, judged only by how badly the plan breaks
if each is false, never by what kind of fault it represents---at three frozen
settings (Table~\ref{tab:v3-overhead}).

\begin{table}[t]
\caption{Live clean-run overhead by checklist-coverage cap (GPT-5.5 writer,
reasoning fixed at \texttt{low}). Batch (S1) clean median \$0.2219 is the
denominator, measured live. Single seed $\times$ 4 tasks (departure D-V3-1); the
writer/batch-worker column is the robust reading (see text).}
\label{tab:v3-overhead}
\begin{tabular}{lrrrr}
\toprule
coverage & assumptions & writer \$ & overhead & overhead \\
 & (median) & (median) & (frozen) & (writer/batch) \\
\midrule
full   & 40 & 0.093 & $+33.9\%$ & $+42.0\%$ \\
cap-12 & 12 & 0.049 & $+20.6\%$ & $+21.9\%$ \\
cap-6  &  6 & 0.038 & $+22.4\%$ & $+17.1\%$ \\
\bottomrule
\end{tabular}
\end{table}

\textbf{No setting clears 12\%.} The raw
$(\mathrm{monitored}-\mathrm{batch})/\mathrm{batch}$ median is noisy at this
study's single seed---worker cost varies run to run and swamps the writer
saving, so cap-6 even reads worse than cap-12 on that column---so we read the
overhead the robust way, as the writer call against the batch worker cost it adds
to: $+42\%$, then $+22\%$, then $+17\%$, every setting above the bar. The floor is
the writer's own input tokens; even a six-assumption checklist costs \$0.038.

And the cap that lowers cost costs detection
(Table~\ref{tab:v3-detection}).

\begin{table}[t]
\caption{Live detection---an interrupt fired, attributed to the fault, during
full worker execution---by checklist-coverage cap. Five catchable faults plus the
DEPENDENCY\_VERSION wall.}
\label{tab:v3-detection}
\begin{tabular}{lccc}
\toprule
injected fault & full & cap-12 & cap-6 \\
\midrule
removed endpoint (API\_SURFACE)        & fires & fires & \textbf{miss} \\
expired token (PERMISSION\_AUTH)       & fires & fires & fires \\
skipped gate (TOOL\_CONTRACT)          & fires & fires & fires \\
swapped passage (RETRIEVAL\_INTEGRITY) & fires & \textbf{miss} & \textbf{miss} \\
RESOURCE\_BUDGET (quota cliff)         & fires & fires & fires \\
DEPENDENCY\_VERSION (silent bump)      & miss  & miss  & miss  \\
\midrule
catchable fired                        & 5/5   & 4/5   & 3/5   \\
\bottomrule
\end{tabular}
\end{table}

\textbf{Detection degrades monotonically as coverage drops: 5, then 4, then 3 of
the five catchable faults.} Capping to twelve assumptions drops the swapped
passage---the cap kept value checks elsewhere but not on the one passage that was
altered. Capping to six also drops the removed endpoint---six assumptions no
longer reach the pricing surface at all. The one setting that catches all five is
the full checklist, which has the \emph{worst} overhead. Cheapening the writer
loses faults and \emph{still} misses the bar.

\textbf{Verdict, against the frozen V3 bar: the coupling holds (live).} Both
committed failure conditions are met---detection erodes as cost drops, and no
setting clears 12\%. This is the paper's thesis on a second model family, and it
arrives through a \emph{new mechanism}. On our own system the cost center was the
one reasoning step (\S8); on GPT-5.5 the reasoning is cheap and detection-free,
and the cost lever is instead the checklist's \emph{length}---which is its
coverage. Different mechanism, same coupling: the thing you would shorten to save
money is the thing that does the catching. The whole run cost \$2.86 in
metered spend, under the pre-registered cap.

\textbf{What this is and is not.} It is a fail on the bar, reported as one, not an
escape. GPT-5.5 is genuinely the better writer; it does not clear the cost. Four
cautions survive verbatim. (1)~The study is a single seed (departure D-V3-1); the
robust signals are the writer-cost overhead and the detection \emph{slope}, not
any single cell. (2)~Prompt-caching did not engage this run; a warm cache could
push the six-assumption cost floor toward 6\%, but that was not achieved here and
is future work, not a result---and even a cost pass would not matter, because six
assumptions already miss two of the five faults. (3)~The offline replay and the
live run are different measurements and we do not blend them; the offline optimism
was contradicted live at the load-bearing point, and that contradiction is part
of the finding, not an embarrassment. (4)~The DEPENDENCY\_VERSION wall of \S7
stands for GPT-5.5 too---missed at every setting, constant across arms, and so
excluded from the coverage slope; and RESOURCE\_BUDGET, which reads 8 of 8 in the
coverage-bound offline replay, fires live at every coverage because its family is
broad enough for even six assumptions to cover it. Each is labeled as what it is.

\section{Threats to Validity}
\textbf{These are findings, not laws.} One harness, one mocked world, three
seeds, bars written on medians. The workers are one model family throughout; the
\emph{check-writer}---the coupling's cost center---we now test on two (\S9), where
the coupling reproduces through a different mechanism. Plan shape varies run to
run and we do not control it. Our injected faults approximate, but do not sample, the real
distribution; qualification guarantees each one breaks the baseline, not that it
is representative. We state the observation and cost results as regularities of
this setting, not proofs.

\textbf{The free-looking mock cuts both ways.} It makes the cost verdict
\emph{conservative}---the system fails the budget where looking is free
(\S8)---but it cannot settle real-deployment cost, where probe volume and
latency return. The 55.5\% is a floor, and the one cost lever is also the
capability, so the cost fix carries detection risk. Both are named; neither is
solved.

\textbf{The redesign learned from the pilot's traces.} That is a real leak, and
we mitigate it rather than hide it: fresh fault settings the designers never
saw, sealed held-out types (drawn by an unseeded cryptographic process with a
committed hash) inside the scoring, and a check-writer whose blindness to the
fault list is verifiable by inspection. The held-out types are held out in their
settings, not their \emph{kind}---same hands, same ontology---and the baseline
defeated two of three designs before one qualified, which shows authoring
genuinely wounding faults is itself hard. A held-out-\emph{ontology} arm and a
natural-failure arm remain the main generalization tests for future work.

\textbf{We built the benchmark, the system, and the bars.} The counterweight is
the apparatus: frozen pre-registration, pre-decided consequences, cryptographic
fingerprints, binding outside review with trace-level rulings, a public hash,
and negatives reported at full prominence---including a naive baseline beating
our own design (\S5). The strongest evidence it binds is behavioral: frozen bars
returned a kill verdict three times, across two model families---twice on our own
system and, under a separately and later frozen pre-registration, a third time on
a different vendor's model (\S9)---and the pre-decided consequences ran as written
every time.

\section{Conclusion}
We pre-registered a system we believed in, and the bars we froze killed most of
it. We diagnosed why from fully recorded runs, rebuilt the detector around the
diagnosis---and the bars killed the rebuild too, on cost. Because each verdict
came with its mechanism instead of a shrug, they compose into one finding.

In budget-bounded multi-agent LLM execution, catching a broken plan is limited
by observation, not by the monitor's cleverness: the system that re-looks,
detects, and the plain baselines that kept re-looking saw what the clever one
starved itself out of seeing. And because the one thing that restores
observation---re-reading the plan to write fresh checks---is also the main cost,
detection and cost are coupled. Restoring observation fixes detection and ends
the self-harm. It does not yet pay for itself, and the obvious way to make it
pay does not open, even when looking is free.

And it is not one model's artifact. A separately and later pre-registered live
study put a different vendor's model in the check-writer's seat---a cheaper,
better writer---and the coupling held: the bar failed a third time, now through
the checklist's length rather than its reasoning (\S9).

That coupling, not any single missed bar, is the result. It is what a
plan-driven monitor has to solve next.

\section*{Data Availability}
A replication package---the benchmark world and harness, the frozen
pre-registration with fingerprints, all 195 pilot traces and the deterministic
replay battery, the 172-run confirmatory experiment and its report, the cost
autopsy, the fan-out sizing pilot, the qualification and review records, the
departure log, and the run log---is at [REPLICATION-PACKAGE-URL]. It also includes
the second-family study of \S9: its own, separately and later frozen
pre-registration (dated 2026-07-01), the live priced matrix and per-run cost and
detection records, and the offline replay scores. The held-out parameter values
stay sealed with someone outside the project; their fingerprint is committed
publicly so the file cannot be swapped.

\appendix

\section{Integrity and Departure Log}
Load-bearing departures, in brief; full records in the departure log (Data
Availability). \textbf{D23} removed two conditionally-qualified runs from the
experiment (215 $\rightarrow$ 172 with four arms) and named the \emph{write-side
single-look} residual---an agent's own write can overwrite a change it never
saw. \textbf{D24} fixed the dead-check class (84 checks, instrument-level, with a
regression test). \textbf{D25/D27} fix the held-out scoring outside the project
and freeze the check-writer's example-selection rule before any tuning.
\textbf{D28--D32} record the confirmation, scheduling, opening-reading,
write-handling, and grounding rulings of the rebuild, each committed before its
code and each preserving identical behavior when the flag is off. \textbf{D33}
set the unbuilt rebuilt-judge arm aside. \textbf{D34} characterizes the lone
replay difference (123/124, benign, on the non-detecting arm); no number or the
overall fail changes.

\textbf{The second-family study (\S9) is pre-registered on its own, and later.}
Its frozen pre-registration is dated 2026-07-01---after the original and both its
verdicts---with its own 12\% bar, arms, coverage settings, verdict definitions,
and a hard spend cap, all committed before its data existed and never folded into
the original pre-registration. \textbf{D-V3-1} logs its execution departures: the
live matrix ran at a single seed (four tasks) with a reduced injected set, and
was executed in resumable foreground chunks after the background runner proved
unreliable---execution-only changes that leave the frozen bar, lever, verdict
meanings, and spend cap untouched. The study spent \$2.86 of metered budget under
its cap; no cap was hit.

\section{Fan-out Break-Even Model (a prediction, not a result)}
Cost-positivity holds when
$C + J + p\cdot R < p\cdot(W_{\mathrm{batch}} - W_{\mathrm{sent}})$. Fitting the
model to the first version's runs ($C$~=~\$0.322 check-writing, $J$~=~\$0.048
judge, $R$~=~\$0.257 per replan) gives a \emph{negative} waste gap of
$-$\$0.072 and \textbf{no crossover at any fan-out, with probability 1.00 over
1,000 resamples}. Re-fitting to the redesign gives a positive gap and an
extrapolated crossover near 86, 40, and 25 workers at fault rates 0.10, 0.25,
and 0.50. \textbf{These are model extrapolations---predictions to test, never
measured results.} The sizing pilot (\S8) tested the direction and closed it
negatively.

\begin{thebibliography}{99}
\bibitem{ref1} Execution monitoring survey (Dunlap et al.).
\bibitem{ref2} Davis-Mendelow et al. Assumption-based planning. AAAI 2013.
\bibitem{ref3} Bercher et al. Plan repair. ICAPS 2014.
\bibitem{ref4} Bauer, Leucker, Schallhart. Runtime verification for LTL and TLTL. ACM TOSEM.
\bibitem{ref5} Havelund, Rosu. Synthesizing monitors for safety properties. TACAS 2002.
\bibitem{ref6} Ferrando, Cardoso. RVPLAN: Runtime Verification of Assumptions in Automated Planning. ICAART 2022, Vol. 2, pp. 67--77. DOI 10.5220/0010776500003116.
\bibitem{ref7} Parallelized planning-acting. arXiv:2503.03505.
\bibitem{ref8} Weak-to-strong monitoring. arXiv:2508.19461.
\bibitem{ref9} AgentSpec. ICSE 2026. arXiv:2503.18666.
\bibitem{ref10} Agent-C. arXiv:2512.23738.
\bibitem{ref11} Pro2Guard. arXiv:2508.00500.
\bibitem{ref12} ProbGuard. arXiv:2602.19844.
\bibitem{ref13} LlamaFirewall. arXiv:2505.03574.
\bibitem{ref14} MAST: multi-agent system failure taxonomy. arXiv:2503.13657.
\bibitem{ref15} AgentOps. arXiv:2411.05285.
\bibitem{ref16} Wasted-computation diagnosis in multi-agent systems.
\bibitem{ref17} LumiMAS. AAMAS 2026. arXiv:2508.12412.
\bibitem{ref18} DiLLS. CHI 2026. arXiv:2602.05446.
\bibitem{ref19} Graph Harness. arXiv:2604.11378.
\bibitem{ref20} SHIELDA. arXiv:2508.07935.
\bibitem{ref21} Learning to Interrupt. arXiv:2604.06452.
\bibitem{ref22} LangGraph interrupt documentation.
\bibitem{ref23} Kephart, Chess. The vision of autonomic computing. IEEE Computer, 2003.
\bibitem{ref24} Weyns. Introduction to self-adaptive systems. 2020.
\bibitem{ref25} Nascimento et al. arXiv:2307.06187.
\bibitem{ref26} ACM TAAS roadmap. DOI 10.1145/3686803.
\bibitem{ref27} Anthropic Engineering. How we built our multi-agent research system.
\bibitem{ref28} Mullick, T\"uz\"un. [Prior work by the authors; citation to be completed at submission.]
\bibitem{ref29} Mullick, T\"uz\"un. [Prior work by the authors; citation to be completed at submission.]
\bibitem{ref30} Zep temporal knowledge graphs. arXiv:2501.13956.
\bibitem{ref31} GAIA. arXiv:2311.12983.
\bibitem{ref32} tau-bench. arXiv:2406.12045.
\bibitem{ref34} Ferrando, Cardoso. RVPLAN: a general purpose framework for replanning using runtime verification. VORTEX@ISSTA 2021, pp. 22--25. DOI 10.1145/3464974.3468447.
\bibitem{ref35} Bozzano, Cimatti, Roveri, Tchaltsev. A comprehensive approach to on-board autonomy verification and validation. IJCAI 2011, pp. 2398--2403.
\bibitem{ref36} Bensalem, Havelund, Orlandini. Verification and validation meet planning and scheduling. STTT 16(1):1--12, 2014.
\bibitem{ref37} Ferrando, Cardoso. Runtime monitoring of action specifications for replanning in classical planning. VORTEX 2025.
\bibitem{ref38} OpenAI. GPT-5.5 model reference and API pricing. developers.openai.com/api/docs/models/gpt-5.5 and /api/docs/pricing (model \texttt{gpt-5.5-2026-04-23}; \$5.00/\$0.50/\$30.00 per 1M input/cached/output tokens; verified 2026-07-01).
\end{thebibliography}

\end{document}
